%%%%%%%%%%%%%%%%%%%%%%%%%%%%%%%%%%%%%%%%%%%%%%%%%%%%%%%%%%%%%%%%%%%%%%%%%%%%%%%%%%
\begin{frame}[fragile]\frametitle{Tadasana}
\begin{columns}
    \begin{column}[T]{0.5\linewidth}
      \begin{itemize}
        \item Stand with feet together, arms by sides.
        \item Distribute weight evenly on both feet.
        \item Engage thighs and lift chest.
        \item Extend arms overhead, palms facing each other.
        \item Hold the pose and breathe deeply.
        \item \textbf{Benefits:} Improves posture, strengthens legs, and enhances concentration.
        \item \textbf{Contraindications:} Avoid rising on the toes with acute cardiac problems, varicose veins or vertigo (Common Yoga Protocol).
      \end{itemize}
    \end{column}
    \begin{column}[T]{0.5\linewidth}
        \begin{center}
		        % FIG ycb1-fig-tadasan | now: images/ycb1_tadasan.png | final: books/YCB1/images/new/ycb1_tadasan_new.png | register: books/YCB1/docs/image_register.csv
		        % PROMPT SP-FIG (common text: books/YCB1/docs/image_system_prompt.md) | Used in block ycb1-pd-yogasana-practice, frame 'Tadasana'. Posture: Tadasana. One adult practitioner (CS1) in the final held position of Tadasana, side view, standing, facing right. The whole body in frame with a small margin; head, spine, arms, hands, legs and feet positions exactly as in the notes block for this posture; weight and support points clear. No props unless the notes block names one. No labels. Aspect ratio 4:5 portrait. Check before use: certified teacher confirms the form against the steps in the notes block.
		        % NEW-PIC (uncomment to show the new picture beside the original; comment the original only after review) \includegraphics[width=0.35\linewidth,keepaspectratio]{new/ycb1_tadasan_new.png}
		        \includegraphics[width=0.35\linewidth,keepaspectratio]{ycb1_tadasan}
				
				{\tiny (Ref: Certification  of Yoga Professionals Official Guidebook For Level I (Instructor))}		
        \end{center}    
    \end{column}
  \end{columns}
\end{frame}

%%%%%%%%%%%%%%%%%%%%%%%%%%%%%%%%%%%%%%%%%%%%%%%%%%%%%%%%%%%%%%%%%%%%%%%%%%%%%%%%%%
\begin{frame}[fragile]\frametitle{Vrikshasana}
\begin{columns}
    \begin{column}[T]{0.5\linewidth}
      \begin{itemize}
        \item Stand in Tadasana position.
        \item Shift weight to one foot, bend the other knee.
        \item Place the sole of the bent foot on the inner thigh of the standing leg.
        \item Join hands in front of the chest or extend overhead.
        \item Hold the position, focus on balance.
        \item \textbf{Benefits:} Enhances balance, strengthens legs, and improves concentration.
        \item \textbf{Contraindications:} Avoid with arthritis, vertigo or obesity (Common Yoga Protocol); also with knee or ankle injury.
      \end{itemize}
    \end{column}
    \begin{column}[T]{0.5\linewidth}
        \begin{center}
		        % FIG ycb1-fig-vrikshasan | now: images/ycb1_vrikshasan.png | final: books/YCB1/images/new/ycb1_vrikshasan_new.png | register: books/YCB1/docs/image_register.csv
		        % PROMPT SP-FIG (common text: books/YCB1/docs/image_system_prompt.md) | Used in block ycb1-pd-yogasana-practice, frame 'Vrikshasana'. Posture: Vrikshasana. One adult practitioner (CS1) in the final held position of Vrikshasana, side view, standing, facing right. The whole body in frame with a small margin; head, spine, arms, hands, legs and feet positions exactly as in the notes block for this posture; weight and support points clear. No props unless the notes block names one. No labels. Aspect ratio 4:5 portrait. Check before use: certified teacher confirms the form against the steps in the notes block.
		        % NEW-PIC (uncomment to show the new picture beside the original; comment the original only after review) \includegraphics[width=0.35\linewidth,keepaspectratio]{new/ycb1_vrikshasan_new.png}
		        \includegraphics[width=0.4\linewidth,keepaspectratio]{ycb1_vrikshasan}
				
				{\tiny (Ref: Certification  of Yoga Professionals Official Guidebook For Level I (Instructor))}	        
		\end{center}    
    \end{column}
  \end{columns}
\end{frame}

%%%%%%%%%%%%%%%%%%%%%%%%%%%%%%%%%%%%%%%%%%%%%%%%%%%%%%%%%%%%%%%%%%%%%%%%%%%%%%%%%%
\begin{frame}[fragile]\frametitle{Ardha Chakrasana}
\begin{columns}
    \begin{column}[T]{0.5\linewidth}
      \begin{itemize}
        \item Stand with feet shoulder-width apart.
        \item Place hands on lower back for support.
        \item Drop the head back; inhale and bend back from the lower back.
        \item Exhale and relax, then come up slowly with an inhalation.
        \item Hold for 10 to 30 seconds with normal breathing.
        \item \textbf{Benefits:} Makes the spine flexible and strengthens the spinal nerves and neck muscles.
        \item \textbf{Contraindications:} Avoid with vertigo or giddiness; those with high blood pressure bend with care (Common Yoga Protocol).
      \end{itemize}
    \end{column}
    \begin{column}[T]{0.5\linewidth}
        \begin{center}
        \begin{center}
		        % FIG ycb1-fig-ardhachakrasan | now: images/ycb1_ardhachakrasan.png | final: books/YCB1/images/new/ycb1_ardhachakrasan_new.png | register: books/YCB1/docs/image_register.csv
		        % PROMPT SP-FIG (common text: books/YCB1/docs/image_system_prompt.md) | Used in block ycb1-pd-yogasana-practice, frame 'Ardha Chakrasana or Hastottanasan'. Posture: Ardha Chakrasana or Hastottanasan. One adult practitioner (CS1) in the final held position of Ardha Chakrasana or Hastottanasan, side view, standing, facing right. The whole body in frame with a small margin; head, spine, arms, hands, legs and feet positions exactly as in the notes block for this posture; weight and support points clear. No props unless the notes block names one. No labels. Aspect ratio 4:5 portrait. Check before use: certified teacher confirms the form against the steps in the notes block.
		        % NEW-PIC (uncomment to show the new picture beside the original; comment the original only after review) \includegraphics[width=0.35\linewidth,keepaspectratio]{new/ycb1_ardhachakrasan_new.png}
		        \includegraphics[width=\linewidth,keepaspectratio]{ycb1_ardhachakrasan}
				
				{\tiny (Ref: Certification  of Yoga Professionals Official Guidebook For Level I (Instructor))}	        
		\end{center}   
        \end{center}    
    \end{column}
  \end{columns}
\end{frame}

%%%%%%%%%%%%%%%%%%%%%%%%%%%%%%%%%%%%%%%%%%%%%%%%%%%%%%%%%%%%%%%%%%%%%%%%%%%%%%%%%%
\begin{frame}[fragile]\frametitle{Padahastasana}
\begin{columns}
    \begin{column}[T]{0.5\linewidth}
      \begin{itemize}
        \item Stand with feet together, arms by sides.
        \item Inhale and raise arms overhead.
        \item Exhale and bend forward, reaching for the feet.
        \item Keep knees slightly bent if needed.
        \item Hold the pose and breathe deeply.
        \item \textbf{Benefits:} Makes the spine flexible and improves digestion (Common Yoga Protocol).
        \item \textbf{Contraindications:} Avoid with cardiac or back problems, hernia, ulcers, abdominal inflammation, high myopia, vertigo or in pregnancy (Common Yoga Protocol).
      \end{itemize}
    \end{column}
    \begin{column}[T]{0.5\linewidth}
        \begin{center}
        \begin{center}
		        % FIG ycb1-fig-padahastasan | now: images/ycb1_padahastasan.png | final: books/YCB1/images/new/ycb1_padahastasan_new.png | register: books/YCB1/docs/image_register.csv
		        % PROMPT SP-FIG (common text: books/YCB1/docs/image_system_prompt.md) | Used in block ycb1-pd-yogasana-practice, frame 'Padahastasana'. Posture: Padahastasana. One adult practitioner (CS1) in the final held position of Padahastasana, side view, standing, facing right. The whole body in frame with a small margin; head, spine, arms, hands, legs and feet positions exactly as in the notes block for this posture; weight and support points clear. No props unless the notes block names one. No labels. Aspect ratio 4:5 portrait. Check before use: certified teacher confirms the form against the steps in the notes block.
		        % NEW-PIC (uncomment to show the new picture beside the original; comment the original only after review) \includegraphics[width=0.35\linewidth,keepaspectratio]{new/ycb1_padahastasan_new.png}
		        \includegraphics[width=0.9\linewidth,keepaspectratio]{ycb1_padahastasan}
				
				{\tiny (Ref: Certification  of Yoga Professionals Official Guidebook For Level I (Instructor))}	        
		\end{center}   
        \end{center}    
    \end{column}
  \end{columns}
\end{frame}

%%%%%%%%%%%%%%%%%%%%%%%%%%%%%%%%%%%%%%%%%%%%%%%%%%%%%%%%%%%%%%%%%%%%%%%%%%%%%%%%%%
\begin{frame}[fragile,shrink=6]\frametitle{Kati Chakrasana}
\begin{columns}
    \begin{column}[T]{0.5\linewidth}
      \begin{itemize}
        \item Stand with feet about half a metre apart; inhale and raise the arms to shoulder level.
        \item Exhale and twist left: right hand to the left shoulder, left arm round the back to the waist. Look over the left shoulder.
        \item Hold the breath out about 2 seconds, inhale and return; repeat on the other side. Do 5 to 10 rounds, without jerking.
        \item \textbf{Benefits:} Tones the waist, back and hips (Satyananda).
        \item \textbf{Contraindications:} Take care with back or spine problems (general caution; the source states none).
      \end{itemize}
    \end{column}
    \begin{column}[T]{0.5\linewidth}
        \begin{center}
        \begin{center}
		        % FIG ycb1-fig-katichakrasan | now: images/ycb1_katichakrasan.png | final: books/YCB1/images/new/ycb1_katichakrasan_new.png | register: books/YCB1/docs/image_register.csv
		        % PROMPT SP-FIG (common text: books/YCB1/docs/image_system_prompt.md) | Used in block ycb1-pd-yogasana-practice, frame 'Kati Chakrasana'. Posture: Kati Chakrasana. One adult practitioner (CS1) in the final held position of Kati Chakrasana, side view, standing, facing right. The whole body in frame with a small margin; head, spine, arms, hands, legs and feet positions exactly as in the notes block for this posture; weight and support points clear. No props unless the notes block names one. No labels. Aspect ratio 4:5 portrait. Check before use: certified teacher confirms the form against the steps in the notes block.
		        % NEW-PIC (uncomment to show the new picture beside the original; comment the original only after review) \includegraphics[width=0.35\linewidth,keepaspectratio]{new/ycb1_katichakrasan_new.png}
		        \includegraphics[width=0.6\linewidth,keepaspectratio]{ycb1_katichakrasan}
				
				{\tiny (Ref: Prana Yoga)}	        
		\end{center}   
        \end{center}    
    \end{column}
  \end{columns}
\end{frame}

%%%%%%%%%%%%%%%%%%%%%%%%%%%%%%%%%%%%%%%%%%%%%%%%%%%%%%%%%%%%%%%%%%%%%%%%%%%%%%%%%%
\begin{frame}[fragile,shrink=6]\frametitle{Trikonasana}
\begin{columns}
    \begin{column}[T]{0.5\linewidth}
      \begin{itemize}
        \item Stand with feet wide apart, arms extended.
        \item Turn one foot out and the other foot slightly in.
        \item Reach down, placing the hand just behind the foot, or on the ankle, shin or knee if the foot is out of reach.
        \item Extend the other arm upwards, gaze up.
        \item Hold the position, then switch sides.
        \item \textbf{Benefits:} Strengthens calf, thigh and waist muscles and makes the spine flexible.
        \item \textbf{Contraindications:} Avoid with a slipped disc, sciatica or after abdominal surgery; do not overdo the lateral stretch.
      \end{itemize}
    \end{column}
    \begin{column}[T]{0.5\linewidth}
        \begin{center}
        \begin{center}
		        % FIG ycb1-fig-trikonasan | now: images/ycb1_trikonasan.png | final: books/YCB1/images/new/ycb1_trikonasan_new.png | register: books/YCB1/docs/image_register.csv
		        % PROMPT SP-FIG (common text: books/YCB1/docs/image_system_prompt.md) | Used in block ycb1-pd-yogasana-practice, frame 'Trikonasana'. Posture: Trikonasana. One adult practitioner (CS1) in the final held position of Trikonasana, side view, standing, facing right. The whole body in frame with a small margin; head, spine, arms, hands, legs and feet positions exactly as in the notes block for this posture; weight and support points clear. No props unless the notes block names one. No labels. Aspect ratio 4:5 portrait. Check before use: certified teacher confirms the form against the steps in the notes block.
		        % NEW-PIC (uncomment to show the new picture beside the original; comment the original only after review) \includegraphics[width=0.35\linewidth,keepaspectratio]{new/ycb1_trikonasan_new.png}
		        \includegraphics[width=\linewidth,keepaspectratio]{ycb1_trikonasan}
				
				{\tiny (Ref: Certification  of Yoga Professionals Official Guidebook For Level I (Instructor))}	        
		\end{center}   
        \end{center}    
    \end{column}
  \end{columns}
\end{frame}

%%%%%%%%%%%%%%%%%%%%%%%%%%%%%%%%%%%%%%%%%%%%%%%%%%%%%%%%%%%%%%%%%%%%%%%%%%%%%%%%%%
\begin{frame}[fragile]\frametitle{Dandasana}
\begin{columns}
    \begin{column}[T]{0.5\linewidth}
      \begin{itemize}
        \item Sit with both legs stretched out and the heels together.
        \item Palms on the floor beside the buttocks.
        \item Keep spine, neck and head upright; close the eyes.
        \item The starting position for the sitting asanas. Relaxed form: legs a little apart, lean slightly back on the hands.
        \item \textbf{Source:} NIOS (secondary source); no benefits or cautions are stated there.
      \end{itemize}
    \end{column}
    \begin{column}[T]{0.5\linewidth}
        \begin{center}
        \begin{center}
		        % FIG ycb1-fig-dandasan | now: images/ycb1_dandasan.png | final: books/YCB1/images/new/ycb1_dandasan_new.png | register: books/YCB1/docs/image_register.csv
		        % PROMPT SP-FIG (common text: books/YCB1/docs/image_system_prompt.md) | Used in block ycb1-pd-yogasana-practice, frame 'Dandasana'. Posture: Dandasana. One adult practitioner (CS1) in the final held position of Dandasana, front three-quarter view, seated on a plain mat. The whole body in frame with a small margin; head, spine, arms, hands, legs and feet positions exactly as in the notes block for this posture; weight and support points clear. No props unless the notes block names one. No labels. Aspect ratio 4:5 portrait. Check before use: certified teacher confirms the form against the steps in the notes block.
		        % NEW-PIC (uncomment to show the new picture beside the original; comment the original only after review) \includegraphics[width=0.35\linewidth,keepaspectratio]{new/ycb1_dandasan_new.png}
		        \includegraphics[width=\linewidth,keepaspectratio]{ycb1_dandasan}
				
				{\tiny (Ref: Prana Yoga)}	        
		\end{center}   
        \end{center}    
    \end{column}
  \end{columns}
\end{frame}

%%%%%%%%%%%%%%%%%%%%%%%%%%%%%%%%%%%%%%%%%%%%%%%%%%%%%%%%%%%%%%%%%%%%%%%%%%%%%%%%%%
\begin{frame}[fragile]\frametitle{Sukhasana}
\begin{columns}
    \begin{column}[T]{0.5\linewidth}
      \begin{itemize}
        \item Sit with legs crossed comfortably.
        \item Place hands on knees or in a mudra.
        \item Keep spine upright and shoulders relaxed.
        \item Close eyes and focus on breath.
        \item Hold the position, breathing deeply.
        \item \textbf{Benefits:} Promotes relaxation, improves flexibility, and calms the mind.
        \item \textbf{Contraindications:} Avoid if you have knee or hip injuries.
      \end{itemize}
    \end{column}
    \begin{column}[T]{0.5\linewidth}
        \begin{center}
        \begin{center}
		        % FIG ycb1-fig-sukhasan | now: images/ycb1_sukhasan.png | final: books/YCB1/images/new/ycb1_sukhasan_new.png | register: books/YCB1/docs/image_register.csv
		        % PROMPT SP-FIG (common text: books/YCB1/docs/image_system_prompt.md) | Used in block ycb1-pd-yogasana-practice, frame 'Sukhasana'. Posture: Sukhasana. One adult practitioner (CS1) in the final held position of Sukhasana, front three-quarter view, seated on a plain mat. The whole body in frame with a small margin; head, spine, arms, hands, legs and feet positions exactly as in the notes block for this posture; weight and support points clear. No props unless the notes block names one. No labels. Aspect ratio 4:5 portrait. Check before use: certified teacher confirms the form against the steps in the notes block.
		        % NEW-PIC (uncomment to show the new picture beside the original; comment the original only after review) \includegraphics[width=0.35\linewidth,keepaspectratio]{new/ycb1_sukhasan_new.png}
		        \includegraphics[width=\linewidth,keepaspectratio]{ycb1_sukhasan}
				
				{\tiny (Ref: Certification  of Yoga Professionals Official Guidebook For Level I (Instructor))}	        
		\end{center}   
        \end{center}    
    \end{column}
  \end{columns}
\end{frame}

%%%%%%%%%%%%%%%%%%%%%%%%%%%%%%%%%%%%%%%%%%%%%%%%%%%%%%%%%%%%%%%%%%%%%%%%%%%%%%%%%%
\begin{frame}[fragile]\frametitle{Padmasana}
\begin{columns}
    \begin{column}[T]{0.5\linewidth}
      \begin{itemize}
        \item Sit with legs extended, then bend one knee.
        \item Place the foot on the opposite thigh.
        \item Repeat with the other leg, placing the foot on the opposite thigh.
        \item Keep spine straight and shoulders relaxed.
        \item Hold the position, focusing on breath.
        \item \textbf{Benefits:} Used for meditation and pranayama; if full lotus is not possible, use half lotus or sukhasana.
        \item \textbf{Contraindications:} Avoid with a knee or hip injury; never force the knees.
      \end{itemize}
    \end{column}
    \begin{column}[T]{0.5\linewidth}
        \begin{center}
        \begin{center}
		        % FIG ycb1-fig-padmasan | now: images/ycb1_padmasan.png | final: books/YCB1/images/new/ycb1_padmasan_new.png | register: books/YCB1/docs/image_register.csv
		        % PROMPT SP-FIG (common text: books/YCB1/docs/image_system_prompt.md) | Used in block ycb1-pd-yogasana-practice, frame 'Padmasana'. Posture: Padmasana. One adult practitioner (CS1) in the final held position of Padmasana, front three-quarter view, seated on a plain mat. The whole body in frame with a small margin; head, spine, arms, hands, legs and feet positions exactly as in the notes block for this posture; weight and support points clear. No props unless the notes block names one. No labels. Aspect ratio 4:5 portrait. Check before use: certified teacher confirms the form against the steps in the notes block.
		        % NEW-PIC (uncomment to show the new picture beside the original; comment the original only after review) \includegraphics[width=0.35\linewidth,keepaspectratio]{new/ycb1_padmasan_new.png}
		        \includegraphics[width=\linewidth,keepaspectratio]{ycb1_padmasan}
				
				{\tiny (Ref: Certification  of Yoga Professionals Official Guidebook For Level I (Instructor))}	        
		\end{center}   
        \end{center}    
    \end{column}
  \end{columns}
\end{frame}

%%%%%%%%%%%%%%%%%%%%%%%%%%%%%%%%%%%%%%%%%%%%%%%%%%%%%%%%%%%%%%%%%%%%%%%%%%%%%%%%%%
\begin{frame}[fragile,shrink=6]\frametitle{Vajrasana}
\begin{columns}
    \begin{column}[T]{0.5\linewidth}
      \begin{itemize}
        \item Kneel on the floor, sit back on heels.
        \item Keep thighs perpendicular to the floor.
        \item Place hands on knees, palms facing down.
        \item Keep spine straight and shoulders relaxed.
        \item Breathe deeply, holding the position.
        \item \textbf{Benefits:} A steady kneeling posture and the starting position for ardha ushtrasana and shashankasana (Common Yoga Protocol).
        \item \textbf{Contraindications:} Avoid with knee or ankle injuries; people with knee osteoarthritis take care or avoid it (Common Yoga Protocol).
      \end{itemize}
    \end{column}
    \begin{column}[T]{0.5\linewidth}
        \begin{center}
        \begin{center}
		        % FIG ycb1-fig-vajrasan | now: images/ycb1_vajrasan.png | final: books/YCB1/images/new/ycb1_vajrasan_new.png | register: books/YCB1/docs/image_register.csv
		        % PROMPT SP-FIG (common text: books/YCB1/docs/image_system_prompt.md) | Used in block ycb1-pd-yogasana-practice, frame 'Vajrasana'. Posture: Vajrasana. One adult practitioner (CS1) in the final held position of Vajrasana, front three-quarter view, seated on a plain mat. The whole body in frame with a small margin; head, spine, arms, hands, legs and feet positions exactly as in the notes block for this posture; weight and support points clear. No props unless the notes block names one. No labels. Aspect ratio 4:5 portrait. Check before use: certified teacher confirms the form against the steps in the notes block.
		        % NEW-PIC (uncomment to show the new picture beside the original; comment the original only after review) \includegraphics[width=0.35\linewidth,keepaspectratio]{new/ycb1_vajrasan_new.png}
		        \includegraphics[width=0.8\linewidth,keepaspectratio]{ycb1_vajrasan}
				
				{\tiny (Ref: Certification  of Yoga Professionals Official Guidebook For Level I (Instructor))}	        
		\end{center}   
        \end{center}    
    \end{column}
  \end{columns}
\end{frame}

%%%%%%%%%%%%%%%%%%%%%%%%%%%%%%%%%%%%%%%%%%%%%%%%%%%%%%%%%%%%%%%%%%%%%%%%%%%%%%%%%%
\begin{frame}[fragile]\frametitle{Bhadrasana}
\begin{columns}
    \begin{column}[T]{0.5\linewidth}
      \begin{itemize}
        \item Sit with legs extended, then bend knees and bring feet together.
        \item Place feet close to the pelvis, holding toes with hands.
        \item Press the knees gently towards the floor; use a soft cushion under them if the thighs do not reach it.
        \item Keep spine erect and shoulders relaxed.
        \item Hold the pose and breathe deeply.
        \item \textbf{Benefits:} Opens hips, improves flexibility, and calms the mind.
        \item \textbf{Contraindications:} Avoid with severe arthritis or sciatica (Common Yoga Protocol); also knee or hip injury.
      \end{itemize}
    \end{column}
    \begin{column}[T]{0.5\linewidth}
        \begin{center}
        \begin{center}
		        % FIG ycb1-fig-bhadrasan | now: images/ycb1_bhadrasan.png | final: books/YCB1/images/new/ycb1_bhadrasan_new.png | register: books/YCB1/docs/image_register.csv
		        % PROMPT SP-FIG (common text: books/YCB1/docs/image_system_prompt.md) | Used in block ycb1-pd-yogasana-practice, frame 'Bhadrasana'. Posture: Bhadrasana. One adult practitioner (CS1) in the final held position of Bhadrasana, front three-quarter view, seated on a plain mat. The whole body in frame with a small margin; head, spine, arms, hands, legs and feet positions exactly as in the notes block for this posture; weight and support points clear. No props unless the notes block names one. No labels. Aspect ratio 4:5 portrait. Check before use: certified teacher confirms the form against the steps in the notes block.
		        % NEW-PIC (uncomment to show the new picture beside the original; comment the original only after review) \includegraphics[width=0.35\linewidth,keepaspectratio]{new/ycb1_bhadrasan_new.png}
		        \includegraphics[width=\linewidth,keepaspectratio]{ycb1_bhadrasan}
				
				{\tiny (Ref: Patanjali Japan Foundation)}	        
		\end{center}   
        \end{center}    
    \end{column}
  \end{columns}
\end{frame}

%%%%%%%%%%%%%%%%%%%%%%%%%%%%%%%%%%%%%%%%%%%%%%%%%%%%%%%%%%%%%%%%%%%%%%%%%%%%%%%%%%
\begin{frame}[fragile,shrink=6]\frametitle{Mandukasana}
\begin{columns}
    \begin{column}[T]{0.5\linewidth}
      \begin{itemize}
        \item Sit in Vajrasana. Make fists, thumbs tucked inside, and press them against the navel.
        \item Exhale and bend forward, gaze ahead; hold the breath out.
        \item Inhale and come up. Repeat three to four times.
        \item \textbf{Benefits:} Flexibility of thighs and legs.
        \item \textbf{Contraindications:} Avoid with abdominal injury, back pain or knee injury.
        \item Source: CBSE Class XI handbook (secondary source).
      \end{itemize}
    \end{column}
    \begin{column}[T]{0.5\linewidth}
        \begin{center}
        \begin{center}
		        % FIG ycb1-fig-mandukasan | now: images/ycb1_mandukasan.png | final: books/YCB1/images/new/ycb1_mandukasan_new.png | register: books/YCB1/docs/image_register.csv
		        % PROMPT SP-FIG (common text: books/YCB1/docs/image_system_prompt.md) | Used in block ycb1-pd-yogasana-practice, frame 'Mandukasana'. Posture: Mandukasana. One adult practitioner (CS1) in the final held position of Mandukasana, front three-quarter view, seated on a plain mat. The whole body in frame with a small margin; head, spine, arms, hands, legs and feet positions exactly as in the notes block for this posture; weight and support points clear. No props unless the notes block names one. No labels. Aspect ratio 4:5 portrait. Check before use: certified teacher confirms the form against the steps in the notes block.
		        % NEW-PIC (uncomment to show the new picture beside the original; comment the original only after review) \includegraphics[width=0.35\linewidth,keepaspectratio]{new/ycb1_mandukasan_new.png}
		        \includegraphics[width=\linewidth,keepaspectratio]{ycb1_mandukasan}
				
				{\tiny (Ref: Atma Bodh)}	        
		\end{center}   
        \end{center}    
    \end{column}
  \end{columns}
\end{frame}

%%%%%%%%%%%%%%%%%%%%%%%%%%%%%%%%%%%%%%%%%%%%%%%%%%%%%%%%%%%%%%%%%%%%%%%%%%%%%%%%%%
\begin{frame}[fragile]\frametitle{Ushtrasana}
\begin{columns}
    \begin{column}[T]{0.5\linewidth}
      \begin{itemize}
        \item Kneel with knees hip-width apart.
        \item Place hands on lower back for support.
        \item Inhale and lift chest, arching back.
        \item Reach for heels with hands, if possible.
        \item Hold the position, breathing deeply.
        \item \textbf{Benefits:} Stretches the entire front body, opens chest, and improves posture.
        \item \textbf{Contraindications:} Avoid with hernia, abdominal injury, arthritis, vertigo or pregnancy (Common Yoga Protocol); also low blood pressure, insomnia or migraine (Guidebook).
      \end{itemize}
    \end{column}
    \begin{column}[T]{0.5\linewidth}
        \begin{center}
        \begin{center}
		        % FIG ycb1-fig-ushtrasan | now: images/ycb1_ushtrasan.png | final: books/YCB1/images/new/ycb1_ushtrasan_new.png | register: books/YCB1/docs/image_register.csv
		        % PROMPT SP-FIG (common text: books/YCB1/docs/image_system_prompt.md) | Used in block ycb1-pd-yogasana-practice, frame 'Ushtrasana'. Posture: Ushtrasana. One adult practitioner (CS1) in the final held position of Ushtrasana, front three-quarter view, seated on a plain mat. The whole body in frame with a small margin; head, spine, arms, hands, legs and feet positions exactly as in the notes block for this posture; weight and support points clear. No props unless the notes block names one. No labels. Aspect ratio 4:5 portrait. Check before use: certified teacher confirms the form against the steps in the notes block.
		        % NEW-PIC (uncomment to show the new picture beside the original; comment the original only after review) \includegraphics[width=0.35\linewidth,keepaspectratio]{new/ycb1_ushtrasan_new.png}
		        \includegraphics[width=\linewidth,keepaspectratio]{ycb1_ushtrasan}
				
				{\tiny (Ref: Certification  of Yoga Professionals Official Guidebook For Level I (Instructor))}	        
		\end{center}   
        \end{center}    
    \end{column}
  \end{columns}
\end{frame}

%%%%%%%%%%%%%%%%%%%%%%%%%%%%%%%%%%%%%%%%%%%%%%%%%%%%%%%%%%%%%%%%%%%%%%%%%%%%%%%%%%
\begin{frame}[fragile]\frametitle{Shashankasana}
\begin{columns}
    \begin{column}[T]{0.5\linewidth}
      \begin{itemize}
        \item Sit in Vajrasana, then spread the knees wide apart with the big toes touching.
        \item Extend arms forward on the floor.
        \item Exhale, bend forward and rest the head on the ground.
        \item Hold the position, breathing deeply.
        \item \textbf{Benefits:} Relieves stress, stretches back and thighs.
        \item \textbf{Contraindications:} Avoid with acute backache; take care with knee osteoarthritis (Common Yoga Protocol).
      \end{itemize}
    \end{column}
    \begin{column}[T]{0.5\linewidth}
        \begin{center}
        \begin{center}
		        % FIG ycb1-fig-shashankasan | now: images/ycb1_shashankasan.png | final: books/YCB1/images/new/ycb1_shashankasan_new.png | register: books/YCB1/docs/image_register.csv
		        % PROMPT SP-FIG (common text: books/YCB1/docs/image_system_prompt.md) | Used in block ycb1-pd-yogasana-practice, frame 'Shashankasana'. Posture: Shashankasana. One adult practitioner (CS1) in the final held position of Shashankasana, front three-quarter view, seated on a plain mat. The whole body in frame with a small margin; head, spine, arms, hands, legs and feet positions exactly as in the notes block for this posture; weight and support points clear. No props unless the notes block names one. No labels. Aspect ratio 4:5 portrait. Check before use: certified teacher confirms the form against the steps in the notes block.
		        % NEW-PIC (uncomment to show the new picture beside the original; comment the original only after review) \includegraphics[width=0.35\linewidth,keepaspectratio]{new/ycb1_shashankasan_new.png}
		        \includegraphics[width=\linewidth,keepaspectratio]{ycb1_shashankasan}
				
				{\tiny (Ref: Certification  of Yoga Professionals Official Guidebook For Level I (Instructor))}	        
		\end{center}   
        \end{center}    
    \end{column}
  \end{columns}
\end{frame}

%%%%%%%%%%%%%%%%%%%%%%%%%%%%%%%%%%%%%%%%%%%%%%%%%%%%%%%%%%%%%%%%%%%%%%%%%%%%%%%%%%
\begin{frame}[fragile]\frametitle{Uttana Mandukasana}
\begin{columns}
    \begin{column}[T]{0.5\linewidth}
      \begin{itemize}
        \item Upright frog: a sitting posture in Vajrasana with the knees apart and the arms folded behind the head.
        \item Steps: source needed (to be confirmed by a certified teacher).
        \item Not the forward bend with the forehead on the floor.
      \end{itemize}
    \end{column}
    \begin{column}[T]{0.5\linewidth}
        \begin{center}
        \begin{center}
		        % FIG ycb1-fig-uttanamandukasan | now: images/ycb1_uttanamandukasan.jpg | final: books/YCB1/images/new/ycb1_uttanamandukasan_new.png | register: books/YCB1/docs/image_register.csv
		        % PROMPT SP-FIG (common text: books/YCB1/docs/image_system_prompt.md) | Used in block ycb1-pd-yogasana-practice, frame 'Uttana Mandukasana'. Posture: Uttana Mandukasana. One adult practitioner (CS1) in the final held position of Uttana Mandukasana, front three-quarter view, seated on a plain mat. The whole body in frame with a small margin; head, spine, arms, hands, legs and feet positions exactly as in the notes block for this posture; weight and support points clear. No props unless the notes block names one. No labels. Aspect ratio 4:5 portrait. Check before use: certified teacher confirms the form against the steps in the notes block.
		        % NEW-PIC (uncomment to show the new picture beside the original; comment the original only after review) \includegraphics[width=0.35\linewidth,keepaspectratio]{new/ycb1_uttanamandukasan_new.png}
		        \includegraphics[width=0.9\linewidth,keepaspectratio]{ycb1_uttanamandukasan}
				
				{\tiny (Ref: Certification  of Yoga Professionals Official Guidebook For Level I (Instructor))}	        
		\end{center}   
        \end{center}    
    \end{column}
  \end{columns}
\end{frame}

%%%%%%%%%%%%%%%%%%%%%%%%%%%%%%%%%%%%%%%%%%%%%%%%%%%%%%%%%%%%%%%%%%%%%%%%%%%%%%%%%%
\begin{frame}[fragile]\frametitle{Paschimottanasana}
\begin{columns}
    \begin{column}[T]{0.5\linewidth}
      \begin{itemize}
        \item Sit with legs extended, feet flexed.
        \item Inhale and lengthen spine.
        \item Exhale and bend forward, reaching for feet.
        \item Hold the pose and breathe deeply.
        \item \textbf{Benefits:} Stretches the spine and hamstrings, calms the mind.
        \item \textbf{Contraindications:} Avoid if you have back or hamstring injuries.
      \end{itemize}
    \end{column}
    \begin{column}[T]{0.5\linewidth}
        \begin{center}
        \begin{center}
		        % FIG ycb1-fig-pashchimottansasan | now: images/ycb1_pashchimottansasan.png | final: books/YCB1/images/new/ycb1_pashchimottansasan_new.png | register: books/YCB1/docs/image_register.csv
		        % PROMPT SP-FIG (common text: books/YCB1/docs/image_system_prompt.md) | Used in block ycb1-pd-yogasana-practice, frame 'Paschimottanasana'. Posture: Paschimottanasana. One adult practitioner (CS1) in the final held position of Paschimottanasana, front three-quarter view, seated on a plain mat. The whole body in frame with a small margin; head, spine, arms, hands, legs and feet positions exactly as in the notes block for this posture; weight and support points clear. No props unless the notes block names one. No labels. Aspect ratio 4:5 portrait. Check before use: certified teacher confirms the form against the steps in the notes block.
		        % NEW-PIC (uncomment to show the new picture beside the original; comment the original only after review) \includegraphics[width=0.35\linewidth,keepaspectratio]{new/ycb1_pashchimottansasan_new.png}
		        \includegraphics[width=\linewidth,keepaspectratio]{ycb1_pashchimottansasan}
				
				{\tiny (Ref: Certification  of Yoga Professionals Official Guidebook For Level I (Instructor))}	        
		\end{center}   
        \end{center}    
    \end{column}
  \end{columns}
\end{frame}

%%%%%%%%%%%%%%%%%%%%%%%%%%%%%%%%%%%%%%%%%%%%%%%%%%%%%%%%%%%%%%%%%%%%%%%%%%%%%%%%%%
\begin{frame}[fragile,shrink=6]\frametitle{Purvottanasana}
\begin{columns}
    \begin{column}[T]{0.5\linewidth}
      \begin{itemize}
        \item Sit with legs out, feet together; palms on the floor beside the hips, arms straight.
        \item Lean back and lift the hips into one straight line, feet on the floor, head tilted back.
        \item Breathe in lifting, out lowering; steady breath in the hold.
        \item Sources disagree on the direction of the fingers (forward or back); not settled here.
        \item \textbf{Benefits:} Strengthens wrists, arms, shoulders and back; stretches legs and hips.
        \item \textbf{Contraindications:} Avoid with wrist or neck injury.
        \item Source: Art of Living (secondary source).
      \end{itemize}
    \end{column}
    \begin{column}[T]{0.5\linewidth}
        \begin{center}
        \begin{center}
		        % FIG ycb1-fig-purvottansasan | now: images/ycb1_purvottansasan.png | final: books/YCB1/images/new/ycb1_purvottansasan_new.png | register: books/YCB1/docs/image_register.csv
		        % PROMPT SP-FIG (common text: books/YCB1/docs/image_system_prompt.md) | Used in block ycb1-pd-yogasana-practice, frame 'Purvottanasana'. Posture: Purvottanasana. One adult practitioner (CS1) in the final held position of Purvottanasana, front three-quarter view, seated on a plain mat. The whole body in frame with a small margin; head, spine, arms, hands, legs and feet positions exactly as in the notes block for this posture; weight and support points clear. No props unless the notes block names one. No labels. Aspect ratio 4:5 portrait. Check before use: certified teacher confirms the form against the steps in the notes block.
		        % NEW-PIC (uncomment to show the new picture beside the original; comment the original only after review) \includegraphics[width=0.35\linewidth,keepaspectratio]{new/ycb1_purvottansasan_new.png}
		        \includegraphics[width=\linewidth,keepaspectratio]{ycb1_purvottansasan}
				
				{\tiny (Ref: Ekhart Yoga)}	        
		\end{center}   
        \end{center}    
    \end{column}
  \end{columns}
\end{frame}

%%%%%%%%%%%%%%%%%%%%%%%%%%%%%%%%%%%%%%%%%%%%%%%%%%%%%%%%%%%%%%%%%%%%%%%%%%%%%%%%%%
\begin{frame}[fragile,shrink=6]\frametitle{Vakrasana}
\begin{columns}
    \begin{column}[T]{0.5\linewidth}
      \begin{itemize}
        \item Sit with legs extended and back straight.
        \item Bend one leg and place its foot beside the opposite knee.
        \item Exhale and twist towards the bent knee, bringing the opposite arm around the knee to hold the big toe or place the palm beside the foot; the other palm on the floor behind.
        \item Hold for 10 to 30 seconds with normal breathing, then switch sides.
        \item \textbf{Benefits:} Increases spinal flexibility; the Common Yoga Protocol also mentions help with constipation and dyspepsia.
        \item \textbf{Contraindications:} Avoid with severe back pain, vertebral and disc disorders, after abdominal surgery and during menstruation.
      \end{itemize}
    \end{column}
    \begin{column}[T]{0.5\linewidth}
        \begin{center}
        \begin{center}
		        % FIG ycb1-fig-vakrasan | now: images/ycb1_vakrasan.jpg | final: books/YCB1/images/new/ycb1_vakrasan_new.png | register: books/YCB1/docs/image_register.csv
		        % PROMPT SP-FIG (common text: books/YCB1/docs/image_system_prompt.md) | Used in block ycb1-pd-yogasana-practice, frame 'Vakrasana'. Posture: Vakrasana. One adult practitioner (CS1) in the final held position of Vakrasana, front three-quarter view, seated on a plain mat. The whole body in frame with a small margin; head, spine, arms, hands, legs and feet positions exactly as in the notes block for this posture; weight and support points clear. No props unless the notes block names one. No labels. Aspect ratio 4:5 portrait. Check before use: certified teacher confirms the form against the steps in the notes block.
		        % NEW-PIC (uncomment to show the new picture beside the original; comment the original only after review) \includegraphics[width=0.35\linewidth,keepaspectratio]{new/ycb1_vakrasan_new.png}
		        \includegraphics[width=\linewidth,keepaspectratio]{ycb1_vakrasan}
				
				{\tiny (Ref: CONDE NAST TRAVELLER)}	        
		\end{center}   
        \end{center}    
    \end{column}
  \end{columns}
\end{frame}

%%%%%%%%%%%%%%%%%%%%%%%%%%%%%%%%%%%%%%%%%%%%%%%%%%%%%%%%%%%%%%%%%%%%%%%%%%%%%%%%%%
\begin{frame}[fragile,shrink=6]\frametitle{Gomukhasana}
\begin{columns}
    \begin{column}[T]{0.5\linewidth}
      \begin{itemize}
        \item Sit with the right knee directly above the left.
        \item Left arm behind the back, right arm over the shoulder; try to clasp the fingers.
        \item Raised elbow behind the head, spine upright, breathe normally.
        \item Hold up to 2 minutes (Satyananda), 1 minute (NIOS); then swap sides.
        \item \textbf{Benefits:} Loosens shoulders and neck, opens the chest, improves posture.
        \item \textbf{Contraindications:} Avoid with shoulder or knee injury; with bleeding piles.
      \end{itemize}
    \end{column}
    \begin{column}[T]{0.5\linewidth}
        \begin{center}
        \begin{center}
		        % FIG ycb1-fig-gomukhasan | now: images/ycb1_gomukhasan.jpg | final: books/YCB1/images/new/ycb1_gomukhasan_new.png | register: books/YCB1/docs/image_register.csv
		        % PROMPT SP-FIG (common text: books/YCB1/docs/image_system_prompt.md) | Used in block ycb1-pd-yogasana-practice, frame 'Gomukhasana'. Posture: Gomukhasana. One adult practitioner (CS1) in the final held position of Gomukhasana, front three-quarter view, seated on a plain mat. The whole body in frame with a small margin; head, spine, arms, hands, legs and feet positions exactly as in the notes block for this posture; weight and support points clear. No props unless the notes block names one. No labels. Aspect ratio 4:5 portrait. Check before use: certified teacher confirms the form against the steps in the notes block.
		        % NEW-PIC (uncomment to show the new picture beside the original; comment the original only after review) \includegraphics[width=0.35\linewidth,keepaspectratio]{new/ycb1_gomukhasan_new.png}
		        \includegraphics[width=\linewidth,keepaspectratio]{ycb1_gomukhasan}
				
				{\tiny (Ref: Himalayan Yoga Ashram)}	        
		\end{center}   
        \end{center}    
    \end{column}
  \end{columns}
\end{frame}

%%%%%%%%%%%%%%%%%%%%%%%%%%%%%%%%%%%%%%%%%%%%%%%%%%%%%%%%%%%%%%%%%%%%%%%%%%%%%%%%%%
\begin{frame}[fragile]\frametitle{Bhujangasana}
\begin{columns}
    \begin{column}[T]{0.5\linewidth}
      \begin{itemize}
        \item Lie on your stomach, legs extended, and feet together.
        \item Place hands under shoulders, elbows close to the body.
        \item Inhale and slowly lift the chin and chest up to the navel region.
        \item Hold the pose and breathe deeply.
        \item \textbf{Benefits:} Helps with stress management and constipation; strengthens the back (Common Yoga Protocol).
        \item \textbf{Contraindications:} Avoid for two to three months after abdominal surgery, with hernia or ulcers, or with back or wrist injury.
      \end{itemize}
    \end{column}
    \begin{column}[T]{0.5\linewidth}
        \begin{center}
        \begin{center}
		        % FIG ycb1-fig-bhujangasan | now: images/ycb1_bhujangasan.png | final: books/YCB1/images/new/ycb1_bhujangasan_new.png | register: books/YCB1/docs/image_register.csv
		        % PROMPT SP-FIG (common text: books/YCB1/docs/image_system_prompt.md) | Used in block ycb1-pd-yogasana-practice, frame 'Bhujangasana'. Posture: Bhujangasana. One adult practitioner (CS1) in the final held position of Bhujangasana, side view, lying face down on a plain mat. The whole body in frame with a small margin; head, spine, arms, hands, legs and feet positions exactly as in the notes block for this posture; weight and support points clear. No props unless the notes block names one. No labels. Aspect ratio 3:2 landscape. Check before use: certified teacher confirms the form against the steps in the notes block.
		        % NEW-PIC (uncomment to show the new picture beside the original; comment the original only after review) \includegraphics[width=0.35\linewidth,keepaspectratio]{new/ycb1_bhujangasan_new.png}
		        \includegraphics[width=\linewidth,keepaspectratio]{ycb1_bhujangasan}
				
				{\tiny (Ref: Certification  of Yoga Professionals Official Guidebook For Level I (Instructor))}	        
		\end{center}   
        \end{center}    
    \end{column}
  \end{columns}
\end{frame}

%%%%%%%%%%%%%%%%%%%%%%%%%%%%%%%%%%%%%%%%%%%%%%%%%%%%%%%%%%%%%%%%%%%%%%%%%%%%%%%%%%
\begin{frame}[fragile]\frametitle{Shalabhasana}
\begin{columns}
    \begin{column}[T]{0.5\linewidth}
      \begin{itemize}
        \item Lie on your stomach in Makarasana, chin on the floor, hands beside the body with palms up.
        \item Inhale and raise the legs off the floor without bending the knees.
        \item Squeeze the buttocks and pull up the knee caps; stay for 10 to 20 seconds.
        \item Hold the position, breathing deeply.
        \item \textbf{Benefits:} Strengthens lower back, improves posture.
        \item \textbf{Contraindications:} Avoid with cardiac problems, high blood pressure, peptic ulcers or hernia; take care with severe lower back pain.
      \end{itemize}
    \end{column}
    \begin{column}[T]{0.5\linewidth}
        \begin{center}
        \begin{center}
		        % FIG ycb1-fig-shalabhasan | now: images/ycb1_shalabhasan.png | final: books/YCB1/images/new/ycb1_shalabhasan_new.png | register: books/YCB1/docs/image_register.csv
		        % PROMPT SP-FIG (common text: books/YCB1/docs/image_system_prompt.md) | Used in block ycb1-pd-yogasana-practice, frame 'Shalabhasana'. Posture: Shalabhasana. One adult practitioner (CS1) in the final held position of Shalabhasana, side view, lying face down on a plain mat. The whole body in frame with a small margin; head, spine, arms, hands, legs and feet positions exactly as in the notes block for this posture; weight and support points clear. No props unless the notes block names one. No labels. Aspect ratio 3:2 landscape. Check before use: certified teacher confirms the form against the steps in the notes block.
		        % NEW-PIC (uncomment to show the new picture beside the original; comment the original only after review) \includegraphics[width=0.35\linewidth,keepaspectratio]{new/ycb1_shalabhasan_new.png}
		        \includegraphics[width=\linewidth,keepaspectratio]{ycb1_shalabhasan}
				
				{\tiny (Ref: Certification  of Yoga Professionals Official Guidebook For Level I (Instructor))}	        
		\end{center}   
        \end{center}    
    \end{column}
  \end{columns}
\end{frame}

%%%%%%%%%%%%%%%%%%%%%%%%%%%%%%%%%%%%%%%%%%%%%%%%%%%%%%%%%%%%%%%%%%%%%%%%%%%%%%%%%%
\begin{frame}[fragile]\frametitle{Makarasana}
\begin{columns}
    \begin{column}[T]{0.5\linewidth}
      \begin{itemize}
        \item Lie on your stomach with the feet wide apart, toes pointing outward.
        \item Bend both arms and place one hand on the other.
        \item Rest the forehead on the hands and close the eyes.
        \item Breathe deeply and relax.
        \item \textbf{Benefits:} Relaxes the lower back; the relaxation posture for the prone asanas.
        \item \textbf{Contraindications:} Avoid with low blood pressure, severe cardiac problems or pregnancy (Common Yoga Protocol).
      \end{itemize}
    \end{column}
    \begin{column}[T]{0.5\linewidth}
        \begin{center}
        \begin{center}
		        % FIG ycb1-fig-makarasan | now: images/ycb1_makarasan.png | final: books/YCB1/images/new/ycb1_makarasan_new.png | register: books/YCB1/docs/image_register.csv
		        % PROMPT SP-FIG (common text: books/YCB1/docs/image_system_prompt.md) | Used in block ycb1-pd-yogasana-practice, frame 'Makarasana'. Posture: Makarasana. One adult practitioner (CS1) in the final held position of Makarasana, side view, lying face down on a plain mat. The whole body in frame with a small margin; head, spine, arms, hands, legs and feet positions exactly as in the notes block for this posture; weight and support points clear. No props unless the notes block names one. No labels. Aspect ratio 3:2 landscape. Check before use: certified teacher confirms the form against the steps in the notes block.
		        % NEW-PIC (uncomment to show the new picture beside the original; comment the original only after review) \includegraphics[width=0.35\linewidth,keepaspectratio]{new/ycb1_makarasan_new.png}
		        \includegraphics[width=\linewidth,keepaspectratio]{ycb1_makarasan}
				
				{\tiny (Ref: Vydya Health)}	        
		\end{center}   
        \end{center}    
    \end{column}
  \end{columns}
\end{frame}

%%%%%%%%%%%%%%%%%%%%%%%%%%%%%%%%%%%%%%%%%%%%%%%%%%%%%%%%%%%%%%%%%%%%%%%%%%%%%%%%%%
\begin{frame}[fragile]\frametitle{Pavanamuktasana}
\begin{columns}
    \begin{column}[T]{0.5\linewidth}
      \begin{itemize}
        \item Lie on your back, knees bent, and feet on the floor.
        \item Hug knees to chest, interlace fingers around shins.
        \item Exhale and lift the head until the chin or nose nears the knees.
        \item Hold the position and breathe deeply.
        \item \textbf{Benefits:} Helps remove flatulence and constipation and stretches the lumbar region (Common Yoga Protocol).
        \item \textbf{Contraindications:} Avoid with abdominal injuries, hernia, sciatica, severe back pain or in pregnancy.
      \end{itemize}
    \end{column}
    \begin{column}[T]{0.5\linewidth}
        \begin{center}
        \begin{center}
		        % FIG ycb1-fig-pavanamuktasan | now: images/ycb1_pavanamuktasan.png | final: books/YCB1/images/new/ycb1_pavanamuktasan_new.png | register: books/YCB1/docs/image_register.csv
		        % PROMPT SP-FIG (common text: books/YCB1/docs/image_system_prompt.md) | Used in block ycb1-pd-yogasana-practice, frame 'Pavanamuktasana'. Posture: Pavanamuktasana. One adult practitioner (CS1) in the final held position of Pavanamuktasana, side view, lying on the back on a plain mat. The whole body in frame with a small margin; head, spine, arms, hands, legs and feet positions exactly as in the notes block for this posture; weight and support points clear. No props unless the notes block names one. No labels. Aspect ratio 3:2 landscape. Check before use: certified teacher confirms the form against the steps in the notes block.
		        % NEW-PIC (uncomment to show the new picture beside the original; comment the original only after review) \includegraphics[width=0.35\linewidth,keepaspectratio]{new/ycb1_pavanamuktasan_new.png}
		        \includegraphics[width=\linewidth,keepaspectratio]{ycb1_pavanamuktasan}
				
				{\tiny (Ref: Certification  of Yoga Professionals Official Guidebook For Level I (Instructor))}	        
		\end{center}   
        \end{center}    
    \end{column}
  \end{columns}
\end{frame}

%%%%%%%%%%%%%%%%%%%%%%%%%%%%%%%%%%%%%%%%%%%%%%%%%%%%%%%%%%%%%%%%%%%%%%%%%%%%%%%%%%
\begin{frame}[fragile]\frametitle{Uttanapadasana / Ardha Halasana}
\begin{columns}
    \begin{column}[T]{0.5\linewidth}
      \begin{itemize}
        \item Lie on your back, legs extended, arms by the sides. The Board syllabus lists uttanapadasana and ardha halasana as two separate postures; the steps here cover the raised-leg form.
        \item Inhale and raise the legs, knees straight, to a comfortable height.
        \item Keep back and shoulders on the floor.
        \item Hold the position and breathe deeply.
        \item \textbf{Benefits:} Strengthens abdominal muscles, tones legs.
        \item \textbf{Contraindications:} Avoid if you have back or leg issues.
      \end{itemize}
    \end{column}
    \begin{column}[T]{0.5\linewidth}
        \begin{center}
        \begin{center}
		        % FIG ycb1-fig-uttanpadasan | now: images/ycb1_uttanpadasan.jpg | final: books/YCB1/images/new/ycb1_uttanpadasan_new.png | register: books/YCB1/docs/image_register.csv
		        % PROMPT SP-FIG (common text: books/YCB1/docs/image_system_prompt.md) | Used in block ycb1-pd-yogasana-practice, frame 'Uttanapadasana'. Posture: Uttanapadasana. One adult practitioner (CS1) in the final held position of Uttanapadasana, side view, lying on the back on a plain mat. The whole body in frame with a small margin; head, spine, arms, hands, legs and feet positions exactly as in the notes block for this posture; weight and support points clear. No props unless the notes block names one. No labels. Aspect ratio 3:2 landscape. Check before use: certified teacher confirms the form against the steps in the notes block.
		        % NEW-PIC (uncomment to show the new picture beside the original; comment the original only after review) \includegraphics[width=0.35\linewidth,keepaspectratio]{new/ycb1_uttanpadasan_new.png}
		        \includegraphics[width=\linewidth,keepaspectratio]{ycb1_uttanpadasan}
				
				{\tiny (Ref: Bodhi School of Yoga)}	        
		\end{center}   
        \end{center}    
    \end{column}
  \end{columns}
\end{frame}

% %%%%%%%%%%%%%%%%%%%%%%%%%%%%%%%%%%%%%%%%%%%%%%%%%%%%%%%%%%%%%%%%%%%%%%%%%%%%%%%%%%
% \begin{frame}[fragile]\frametitle{}
% \begin{columns}
    % \begin{column}[T]{0.5\linewidth}
      % \begin{itemize}
        % \item Lie on your back, arms by sides.
        % \item Lift legs overhead, bringing toes to the floor behind the head.
        % \item Support lower back with hands if needed.
        % \item Hold the position and breathe deeply.
        % \item \textbf{Benefits:} Stretches back and legs, strengthens core.
        % \item \textbf{Contraindications:} Avoid if you have neck or back issues.
      % \end{itemize}
    % \end{column}
    % \begin{column}[T]{0.5\linewidth}
        % \begin{center}
        % \begin{center}
		        % \includegraphics[width=\linewidth,keepaspectratio]{ycb1_vrikshasan}
				
				% {\tiny (Ref: Certification  of Yoga Professionals Official Guidebook For Level I (Instructor))}	        
		% \end{center}   
        % \end{center}    
    % \end{column}
  % \end{columns}
% \end{frame}

%%%%%%%%%%%%%%%%%%%%%%%%%%%%%%%%%%%%%%%%%%%%%%%%%%%%%%%%%%%%%%%%%%%%%%%%%%%%%%%%%%
\begin{frame}[fragile,shrink=6]\frametitle{Setubandhasana}
\begin{columns}
    \begin{column}[T]{0.5\linewidth}
      \begin{itemize}
        \item Lie on your back, knees bent, heels near the buttocks; hold both ankles.
        \item Lift hips towards the ceiling, pressing into feet.
        \item Keep head and shoulders on the floor; the hands may support the waist (a variant interlaces the fingers under the back).
        \item Hold the position and breathe deeply.
        \item \textbf{Benefits:} Strengthens the lower back and helps digestion and constipation (Common Yoga Protocol).
        \item \textbf{Contraindications:} Avoid with ulcers or hernia and in late pregnancy (Common Yoga Protocol); also neck or back injuries.
      \end{itemize}
    \end{column}
    \begin{column}[T]{0.5\linewidth}
        \begin{center}
        \begin{center}
		        % FIG ycb1-fig-setubandhasan | now: images/ycb1_setubandhasan.jpg | final: books/YCB1/images/new/ycb1_setubandhasan_new.png | register: books/YCB1/docs/image_register.csv
		        % PROMPT SP-FIG (common text: books/YCB1/docs/image_system_prompt.md) | Used in block ycb1-pd-yogasana-practice, frame 'Setubandhasana'. Posture: Setubandhasana. One adult practitioner (CS1) in the final held position of Setubandhasana, side view, lying on the back on a plain mat. The whole body in frame with a small margin; head, spine, arms, hands, legs and feet positions exactly as in the notes block for this posture; weight and support points clear. No props unless the notes block names one. No labels. Aspect ratio 3:2 landscape. Check before use: certified teacher confirms the form against the steps in the notes block.
		        % NEW-PIC (uncomment to show the new picture beside the original; comment the original only after review) \includegraphics[width=0.35\linewidth,keepaspectratio]{new/ycb1_setubandhasan_new.png}
		        \includegraphics[width=\linewidth,keepaspectratio]{ycb1_setubandhasan}
				
				{\tiny (Ref: Rishikesh Yogis Yogashala)}	        
		\end{center}   
        \end{center}    
    \end{column}
  \end{columns}
\end{frame}

%%%%%%%%%%%%%%%%%%%%%%%%%%%%%%%%%%%%%%%%%%%%%%%%%%%%%%%%%%%%%%%%%%%%%%%%%%%%%%%%%%
\begin{frame}[fragile,shrink=6]\frametitle{Vipareetakarani}
\begin{columns}
    \begin{column}[T]{0.5\linewidth}
      \begin{itemize}
        \item The Satyananda hands-on-hips version, not legs-up-the-wall.
        \item Lie flat; raise the legs, lift the hips and cup them with the hands, elbows bent.
        \item Legs vertical, trunk at about 45 degrees; weight on shoulders and elbows, chin off the chest.
        \item Breathe normally. Hold a few seconds at first, building to 3 to 5 minutes.
        \item \textbf{Contraindications:} As for Sarvangasana: neck disc problems, high blood pressure, heart problems, weak eye vessels, enlarged thyroid, liver or spleen; avoid in menstruation and late pregnancy.
      \end{itemize}
    \end{column}
    \begin{column}[T]{0.5\linewidth}
        \begin{center}
        \begin{center}
		        % FIG ycb1-fig-vipartikarni | now: images/ycb1_vipartikarni.png | final: books/YCB1/images/new/ycb1_vipartikarni_new.png | register: books/YCB1/docs/image_register.csv
		        % PROMPT SP-FIG (common text: books/YCB1/docs/image_system_prompt.md) | Used in block ycb1-pd-yogasana-practice, frame 'Vipareetakarani'. Posture: Vipareetakarani. One adult practitioner (CS1) in the final held position of Vipareetakarani, side view, lying on the back on a plain mat. The whole body in frame with a small margin; head, spine, arms, hands, legs and feet positions exactly as in the notes block for this posture; weight and support points clear. No props unless the notes block names one. No labels. Aspect ratio 3:2 landscape. Check before use: certified teacher confirms the form against the steps in the notes block.
		        % NEW-PIC (uncomment to show the new picture beside the original; comment the original only after review) \includegraphics[width=0.35\linewidth,keepaspectratio]{new/ycb1_vipartikarni_new.png}
		        \includegraphics[width=\linewidth,keepaspectratio]{ycb1_vipartikarni}
				
				{\tiny (Ref: Yoga4Lyf)}	        
		\end{center}   
        \end{center}    
    \end{column}
  \end{columns}
\end{frame}

%%%%%%%%%%%%%%%%%%%%%%%%%%%%%%%%%%%%%%%%%%%%%%%%%%%%%%%%%%%%%%%%%%%%%%%%%%%%%%%%%%
\begin{frame}[fragile,shrink=6]\frametitle{Saral Matsyasana}
\begin{columns}
    \begin{column}[T]{0.5\linewidth}
      \begin{itemize}
        \item Easy variation of Matsyasana (Satyananda), legs straight. Name matched by us, to be confirmed by a certified teacher.
        \item Stretch both legs straight and lean back on the arms.
        \item Rest the top of the head on the floor and arch the back; hands on thighs or floor.
        \item Return after a short while.
        \item \textbf{Contraindications:} Avoid with heart disease, ulcer, hernia, back conditions, serious illness or pregnancy; take care with blood pressure or a delicate neck.
      \end{itemize}
    \end{column}
    \begin{column}[T]{0.5\linewidth}
        \begin{center}
        \begin{center}
		        % FIG ycb1-fig-saralmatsyasan | now: images/ycb1_saralmatsyasan.jpg | final: books/YCB1/images/new/ycb1_saralmatsyasan_new.png | register: books/YCB1/docs/image_register.csv
		        % PROMPT SP-FIG (common text: books/YCB1/docs/image_system_prompt.md) | Used in block ycb1-pd-yogasana-practice, frame 'Saral Matsyasana'. Posture: Saral Matsyasana. One adult practitioner (CS1) in the final held position of Saral Matsyasana, side view, lying on the back on a plain mat. The whole body in frame with a small margin; head, spine, arms, hands, legs and feet positions exactly as in the notes block for this posture; weight and support points clear. No props unless the notes block names one. No labels. Aspect ratio 3:2 landscape. Check before use: certified teacher confirms the form against the steps in the notes block.
		        % NEW-PIC (uncomment to show the new picture beside the original; comment the original only after review) \includegraphics[width=0.35\linewidth,keepaspectratio]{new/ycb1_saralmatsyasan_new.png}
		        \includegraphics[width=\linewidth,keepaspectratio]{ycb1_saralmatsyasan}
				
				{\tiny (Ref: Kerala Tourism)}	        
		\end{center}   
        \end{center}    
    \end{column}
  \end{columns}
\end{frame}

